\title{QLoRA: Efficient Finetuning of Quantized LLMs}

\begin{abstract}
We introduce \textsc{QLoRA}, a novel and resource-efficient finetuning technique that enables the adaptation of models with up to 65B parameters using a single 48GB GPU, without any compromise on the performance of standard 16-bit finetuning tasks. \textsc{QLoRA} enables gradient backpropagation through a frozen pretrained language model, quantized to 4 bits, and leverages Low Rank Adapters~(LoRA) for adaptation. The top-performing models obtained with this approach, which we call \textbf{Guanaco}, surpass all previously released open models on the Vicuna benchmark, achieving 99.3\% of ChatGPT's performance after just 24 hours of finetuning on a single GPU. \textsc{QLoRA} incorporates several key advancements for memory conservation while maintaining efficacy: (a) 4-bit NormalFloat (NF4), a novel data type optimally designed for representing normally distributed weights, (b) Double Quantization, a technique that further reduces average memory consumption by quantizing quantization parameters themselves, and (c) Paged Optimizers to efficiently handle transient memory surges. Utilizing \textsc{QLoRA}, we finetuned over 1,000 models and performed an extensive evaluation of instruction following and chatbot capabilities, spanning 8 instruction-following datasets, different model architectures (LLaMA, T5), and large-scale models previously impractical for traditional finetuning methods (such as models with 33B and 65B parameters). Our findings demonstrate that QLoRA-based finetuning on compact, high-quality datasets yields state-of-the-art performance, often outperforming larger prior models. We present a comprehensive assessment of chatbot quality employing both human judgment and GPT-4-based evaluations, confirming the latter as an affordable and viable proxy for human annotation. Additionally, we observe that many existing chatbot benchmarks fail to reliably reflect real performance differences. Through a focused error analysis, we illustrate the failure modes where \textbf{Guanaco} lags behind ChatGPT. We provide open access to all models and accompanying code, including CUDA kernels supporting 4-bit model training.\footnote{\url{https://github.com/artidoro/qlora} and \url{https://github.com/TimDettmers/bitsandbytes}}
\end{abstract}

\section{Introduction}

Adapting large language models (LLMs) through finetuning has proven to be an effective strategy for enhancing model capabilities \citep{min2021metaicl, wei2021finetuned, ouyang2022training, wang2022super, wang2022self, liu2022few}, as well as for incorporating or eliminating specific behaviors as desired \citep{ouyang2022training,askell2021general,bai2022training}. Nevertheless, finetuning models of substantial size is often financially unfeasible; for example, standard 16-bit finetuning of the LLaMA model with 65B parameters~\citep{touvron2023llama} demands in excess of 780 GB of GPU memory. Although state-of-the-art quantization techniques have succeeded in lowering the memory requirements for inference in LLMs \citep{dettmers2022llm, dettmers2022case, frantar2022gptq, xiao2022smoothquant}, these methods are not applicable during the training process and tend to fail in this context \citep{wortsman2023stable}.

In this work, we present the first successful approach for finetuning a 4-bit quantized model without loss in performance. Our approach, \textsc{QLoRA}, introduces an innovative high-precision quantization scheme applied to a pretrained model, reducing it to 4 bits. Subsequently, we incorporate a limited number of trainable Low-rank Adapter weights \citep{hu2021lora} to facilitate effective adaptation.

which are optimized by propagating gradients through the quantized weights via backpropagation.

\textsc{QLoRA} significantly lowers the memory demands for finetuning a 65B parameter model, reducing the requirement from over 780GB to under 48GB of GPU memory, all while maintaining both runtime efficiency and predictive accuracy relative to a fully finetuned 16-bit baseline. This advancement transforms the feasibility of LLM finetuning, enabling even the largest publicly accessible models to be finetuned on a single GPU. Utilizing \textsc{QLoRA}, we develop the \textbf{Guanaco} series of models, with the second-best variant attaining 97.8\% of ChatGPT’s performance on the Vicuna~\citep{vicuna2023} benchmark, training in under 12 hours on a standard consumer GPU; with a professional GPU, our largest model achieves 99.3\% within 24 hours, effectively matching ChatGPT’s results on the Vicuna benchmark. In practical deployment, the smallest \textbf{Guanaco} model (7B parameters) requires only 5GB of memory and surpasses a 26GB Alpaca model by more than 20 percentage points on the Vicuna benchmark (see Table~\ref{tbl:vicuna_eval}).

\textsc{QLoRA} incorporates several novel techniques aimed at minimizing memory consumption while maintaining model performance: (1) \textbf{4-bit NormalFloat}, a quantization scheme optimized for data following a normal distribution, shown both theoretically and empirically to outperform 4-bit Integer and 4-bit Float representations. (2) \textbf{Double Quantization}, a strategy in which the quantization coefficients themselves are quantized, achieving an average memory savings of approximately 0.37 bits per parameter (equivalent to about 3 GB for a 65B parameter model). (3) \textbf{Paged Optimizers}, leveraging NVIDIA’s unified memory technology to mitigate the memory surges during gradient checkpointing when handling mini-batches with extended sequence lengths. Collectively, these improvements enable a more finely adjusted LoRA method, placing adapters within every layer of the network, and significantly reducing the loss in accuracy that was common in previous methods.

The efficiency of \textsc{QLoRA} makes it feasible to conduct a comprehensive analysis of instruction finetuning and chatbot capabilities at model sizes that would be prohibitive with standard finetuning methods due to excessive memory requirements. Consequently, we train over 1,000 models spanning multiple instruction tuning corpora, diverse model architectures, and parameter counts ranging from 80M up to 65B. Beyond demonstrating that \textsc{QLoRA} achieves parity with 16-bit precision (\S\ref{sec:comp_to_std}) and facilitates the training of the state-of-the-art chatbot \textbf{Guanaco} (\S\ref{sec:pushing}), we also investigate patterns in the performance of the resulting models. Our findings indicate, firstly, that the quality of the training data significantly outweighs the importance of dataset size; for instance, the OASST1 dataset with 9,000 examples yielded better chatbot outcomes than the FLAN v2 dataset, which was subsampled to 450,000 instances—even though both aimed to promote generalized instruction following. Secondly, we provide evidence that excelling on the Massive Multitask Language Understanding (MMLU) benchmark does not guarantee superior results on the Vicuna chatbot benchmark and vice versa, highlighting that the alignment between dataset and target task is more critical than the overall dataset magnitude.

In addition, we conduct a comprehensive assessment of chatbot capabilities, incorporating evaluations from both human annotators and GPT-4. Our methodology employs a tournament-style framework in which chatbots are pitted against one another in head-to-head contests to generate optimal responses to specific prompts. The victor of each contest is determined by either human judges or GPT-4. These outcomes are then synthesized into Elo scores~\citep{elo1967proposed, elo1978rating}, which serve to rank the chatbots according to their demonstrated performance. Our analysis reveals that, for the most part, the rankings generated by GPT-4 and human evaluators are consistent, although notable discrepancies do occur. Consequently, we emphasize that, while automated model-based evaluation presents a cost-effective alternative to human annotation, it remains subject to certain uncertainties.

We complement our chatbot benchmark findings with a qualitative examination of the \textbf{Guanaco} models. This assessment sheds light on both effective and problematic instances that the quantitative evaluations did not reveal.

We provide access to all generated outputs, complete with both human and GPT-4 evaluations, to support ongoing research. Our codebase and CUDA kernels are openly available, and our techniques have been incorporated into the Hugging Face transformers library \citep{wolf2019huggingface}, ensuring broad usability. Additionally, we publish a suite of adapters compatible with 7/13/33/65B models, each finetuned on one of eight distinct instruction-following datasets, resulting in 32 openly released and specialized model variants.

\section{Background}
\paragraph{Block-wise k-bit Quantization} Quantization refers to the transformation of data from a more informative, higher-precision representation to one with reduced precision and information content. Typically, this involves converting values encoded in a larger number of bits to a format with fewer bits; for instance, transitioning from 32-bit floating point numbers to 8-bit integer values. To maximize the representational capacity of the lower-bit data type, input values are usually rescaled by normalizing with respect to their absolute maximum, ensuring that the full dynamic range of the target type is utilized. Commonly, input elements organized as a tensor undergo this scaling. As an illustration, converting a tensor in the 32-bit floating point (FP32) format to an Int8 tensor with the range $[-127, 127]$ entails:

\begin{equation}
    \mathbf{X}^{ \text{Int8}} = \text{round}\left( \frac{127}{{\text{absmax}}(\mathbf{X}^{{\text{FP32}}})}\mathbf{X}^{{\text{FP32}}} \right) = \text{round}(c^{\text{FP32}}\cdot \mathbf{X}^{{\text{FP32}}}),
\end{equation}

where $c$ denotes the {\em quantization constant} or {\em quantization scale}. The reverse operation, dequantization, is defined as follows:

\begin{equation}
    \text{dequant}(c^{\text{FP32}}, \mathbf{X}^{\text{Int8}}) = \frac{\mathbf{X}^{\text{Int8}}}{c^{\text{FP32}}} = \mathbf{X}^{\text{FP32}}
\end{equation}

A significant limitation of this method is that the presence of a high-magnitude value (an outlier) in the input tensor can result in inefficient utilization of the quantization bins; certain bit patterns may correspond to few or no quantized values. To address the challenge posed by outliers, a widely used strategy is to partition the input tensor into smaller blocks, with each block quantized independently using its own scaling factor $c$. This procedure can be expressed as follows: the input tensor $\mathbf{X}\in \mathbb{R}^{b\times h}$ is divided into $n$ contiguous blocks, each of length $B$, by first flattening the tensor and then slicing it into ${n = ({b\times h} )/{B} }$ linear segments. Each block is then quantized separately, as outlined in Equation~1, resulting in a quantized tensor accompanied by $n$ unique quantization constants $c_i$.

\paragraph{Low-rank Adapters} Low-rank Adapter (LoRA) finetuning \citep{hu2021lora} achieves lower memory usage by introducing a small number of trainable adapter parameters, while the primary model parameters are kept constant and are not updated. During stochastic gradient descent, gradients are backpropagated through the unaltered pretrained weights to the adapter module, allowing only the adapter to be modified in order to minimize the loss. LoRA enhances a standard linear transformation by incorporating an extra, factorized projection. Specifically, for a projection ${\mathbf{X} \mathbf{W} = \mathbf{Y}}$ where $\mathbf{X}\in  \mathbb{R}^{b\times h}$ and $\mathbf{W}\in  \mathbb{R}^{h\times o}$, LoRA modifies the computation as follows:

\begin{equation}
    \mathbf{Y} = \mathbf{X}\mathbf{W} + s\mathbf{X}\mathbf{L}_1\mathbf{L}_2,
\end{equation}

with $\mathbf{L}_1 \in \mathbb{R}^{h \times r}$, $\mathbf{L}_2 \in \mathbb{R}^{r \times o}$, and $s$ representing a scalar value.

\paragraph{Memory Requirement of Parameter-Efficient Finetuning} A key aspect to consider is the memory overhead incurred by LoRA during training, specifically with respect to both the number and capacity of adapters utilized. Due to the exceptionally small memory consumption of LoRA, it is feasible to deploy a greater number of adapters to boost model efficacy while maintaining a modest overall memory demand. Although LoRA was introduced as a Parameter Efficient Finetuning (PEFT) approach, the dominant contributor to memory usage during LLM finetuning is the storage of activation gradients, rather than the LoRA parameters themselves. For instance, with a 7B LLaMA model trained on FLAN v2, using a batch size of 1 and LoRA parameters corresponding to the standard 0.2\% of original model weights~\citep{hu2021lora,liu2022few}, the memory required for storing LoRA input gradients amounts to 567 MB, while the LoRA weights themselves occupy a mere 26 MB. Implementing gradient checkpointing~\citep{chen2016training} brings the per-sequence input gradients down to roughly 18 MB, yet these still exceed the collective memory of all LoRA weights. In comparison, the 4-bit model backbone demands 5,048 MB of memory. These findings illustrate that while gradient checkpointing provides substantial memory savings, further minimizing LoRA parameter size contributes little additional benefit in terms of total memory. Consequently, adopting additional adapters does not significantly impact the training memory footprint (refer to Appendix~\ref{app:memory} for comprehensive figures). As will be discussed subsequently, this flexibility is essential for attaining full 16-bit precision performance.

\section{\textsc{QLoRA} Finetuning}

\textsc{QLoRA} enables precise finetuning with 4-bit representations by leveraging two novel methods we introduce: 4-bit NormalFloat (NF4) quantization and Double Quantization. Furthermore, we present Paged Optimizers as a solution to memory spikes during gradient checkpointing, which have historically led to out-of-memory issues and impeded single-machine finetuning of large-scale models.

\textsc{QLoRA} utilizes a low-precision storage format—typically 4-bit in our implementation—while relying on a computation data type, most often BFloat16. Consequently, each time a \textsc{QLoRA} weight tensor is accessed, it is first dequantized into BFloat16 before carrying out the matrix multiplication at 16-bit precision.

In the following, we examine the individual elements that constitute \textsc{QLoRA}, and subsequently provide a formal definition of \textsc{QLoRA}.

\paragraph{4-bit NormalFloat Quantization} The NormalFloat (NF) data type extends the concept of Quantile Quantization \citep{dettmers2022optimizers}, which represents an information-theoretically optimal scheme by guaranteeing that every quantization bin contains the same quantity of values mapped from the original tensor. This approach operates by leveraging the empirical cumulative distribution function to estimate quantiles within the input tensor.

A significant drawback of quantile quantization lies in the high computational cost associated with quantile estimation. To mitigate this, efficient quantile approximation methods like SRAM quantiles~\citep{dettmers2022optimizers} are commonly employed. However, since these algorithms only provide approximate quantile values, the resulting data types tend to exhibit substantial quantization errors, particularly for outlier values—elements that are frequently of highest significance.

When input tensors originate from a distribution that remains constant aside from a quantization factor, costly quantile computations and related approximation inaccuracies can be circumvented. Under these circumstances, the quantiles of the input tensors align, allowing for precise quantile calculation with manageable computational demands.

Given that pretrained neural network weights typically follow a zero-mean normal distribution characterized by a standard deviation $\sigma$ (refer to Appendix~\ref{app:normality}), we standardize all weights to a unified distribution by adjusting $\sigma$ so the resulting values align precisely within our chosen data type range. Specifically, we define the range for our data type as $[-1, 1]$. Therefore, it is necessary to normalize both the data type quantiles and the neural network weights so they reside within this interval.

The process for determining the data type that achieves information-theoretic optimality for zero-mean normal distributions with arbitrary standard deviation $\sigma$ over the interval $[-1, 1]$ involves several steps: (1) First, calculate the $2^k+1$ quantile points of a standard normal distribution $N(0, 1)$ to produce a $k$-bit quantile-based quantization data type tailored to normal distributions; (2) Next, map the values of this quantized data type onto the interval $[-1, 1]$ by normalizing accordingly; (3) Finally, given an input weight tensor, rescale it to the $[-1, 1]$ interval using absolute maximum normalization prior to quantization.

After aligning both the range of the weights and the data type, standard quantization procedures can be applied. Essentially, step (3) involves adjusting the standard deviation of the weight tensor so that it corresponds to that of the $k$-bit data type. To express this rigorously, we determine the $2^k$ quantization values $q_i$ for the data type using the following approach:

\begin{equation}
q_i  = \frac{1}{2}\left( Q_X\left(\frac{i}{2^k+1}\right) + Q_X\left(\frac{i+1}{2^k+1}\right)\right),
\end{equation}

Here, $Q_X(\cdot)$ denotes the quantile function associated with the standard normal distribution $N(0,1)$. In the context of symmetric k-bit quantization, a notable limitation is the lack of a precise representation for zero—an essential characteristic for accurately quantizing padding or elements inherently valued at zero, without introducing quantization error. To address this, and to guarantee that zero is mapped to a specific discrete quantization point while fully utilizing all $2^k$ possible bit values in the k-bit type, we propose constructing an asymmetric data type. This involves separately computing quantiles $q_i$ for negative values (with $2^{k-1}$ quantiles) and for positive values (with $2^{k-1}+1$ quantiles), then merging these two sets and removing one of the duplicated zeros from the overlap. We refer to the resulting format as the \emph{k-bit NormalFloat} (NFk), as it guarantees that each quantization interval has the same expected frequency under a zero-mean normal distribution and is, from an information-theoretic standpoint, optimal for such data. The specific quantized values of this data type are detailed in Appendix~\ref{app:nf4}.

\paragraph{Double Quantization{}}
\emph{Double Quantization} (DQ) refers to the technique of further quantizing the quantization constants themselves in order to achieve greater memory efficiency. Accurate 4-bit quantization necessitates a relatively small blocksize~\citep{dettmers2022case}, which, however, significantly increases memory consumption. For instance, if 32-bit quantization constants are employed with a blocksize of 64 for $\mathbf{W}$, the storage requirement for these constants averages out to $32/64 = 0.5$ bits per parameter. By applying Double Quantization, it becomes possible to reduce the memory overhead incurred by quantization constants.

In greater detail, Double Quantization involves taking the initial quantization constants $c_2^{\text{FP32}}$ and subjecting them to a subsequent quantization procedure. This secondary quantization produces both the quantized constants $c_2^{\text{FP8}}$ and a new set of quantization constants, $c_1^{\text{FP32}}$. For this secondary quantization phase, we employ 8-bit floating point numbers with a blocksize of 256, as experimental results indicate there is no noticeable drop in performance with 8-bit quantization, corroborating findings from \citet{dettmers2022case}. 

Given that the $c_2^{\text{FP32}}$ values are non-negative, we center them around zero by subtracting their mean prior to quantization, enabling the use of symmetric quantization schemes. For a blocksize of 64, this approach decreases the per-parameter memory requirement from $32/64=0.5$ bits to $8/64 + 32/(64\cdot256) = 0.127$ bits, corresponding to a savings of 0.373 bits for each parameter.

\paragraph{Paged Optimizers} leverage the NVIDIA unified memory~\footnote{\tiny \url{https://docs.nvidia.com/cuda/cuda-c-programming-guide}} capability, which transparently manages page-level data transfers between CPU and GPU. This functionality ensures reliable GPU execution, even when the GPU intermittently exhausts its available memory. Analogous to traditional paging between main memory and disk storage, this mechanism permits the allocation of paged memory for optimizer states. When GPU memory is depleted, these states are seamlessly transferred to CPU RAM and subsequently moved back into GPU memory as required during optimizer updates.

\paragraph{\textsc{QLoRA}.} Building on the previously outlined components, we formalize \textsc{QLoRA} in the context of a quantized base model by specifying its structure for an individual linear layer augmented with a single LoRA adapter as presented below:

\begin{equation}
    \mathbf{Y}^{\text{BF16}} =  \mathbf{X}^{\text{BF16}} \text{doubleDequant}(c_1^{\text{FP32}}, c_2^{\text{k-bit}}, \mathbf{W}^{\text{NF4}}) + \mathbf{X}^{\text{BF16}}\mathbf{L}^{\text{BF16}}_1\mathbf{L}^{\text{BF16}}_2,
\end{equation}

where the function doubleDequant$(\cdot)$ is specified as follows:

\begin{equation}
    \text{doubleDequant}(c_1^{\text{FP32}}, c_2^{\text{k-bit}}, \mathbf{W}^{\text{k-bit}}) = \text{dequant}(\text{dequant}(c_1^{\text{FP32}}, c_2^{\text{k-bit}}), \mathbf{W}^{\text{4bit}}) = \mathbf{W}^{\text{BF16}},
\end{equation}

We represent $\mathbf{W}$ using NF4 and employ FP8 for $c_2$.

A block size of 64 is selected for $\mathbf{W}$ to achieve greater quantization accuracy, whereas a larger block size of 256 is utilized for $c_2$ in order to optimize memory usage.

Parameter updates require computing the gradient of the error only with respect to the adapter weights, $\frac{\partial E}{\partial \mathbf{L}_i}$, rather than the 4-bit weights, $\frac{\partial E}{\partial \mathbf{W}}$. Nonetheless, to obtain $\frac{\partial E}{\partial \mathbf{L}_i}$, it is necessary to also compute $\frac{\partial \mathbf{X}}{\partial \mathbf{W}}$. This process utilizes equation~(5), involving the dequantization of $\mathbf{W}^{\text{NF4}}$ from its storage format to the computation data type $\mathbf{W}^{\text{BF16}}$, so that the derivative $\frac{\partial \mathbf{X}}{\partial \mathbf{W}}$ can be evaluated with BFloat16-level precision.

In summary, \textsc{QLoRA} employs a single data type for storage (commonly 4-bit NormalFloat) alongside a computation data type (16-bit BrainFloat). During both forward and backward propagation, the storage data type is dequantized to the computation data type; however, weight gradients are calculated exclusively for the LoRA parameters, which are represented in 16-bit BrainFloat.

\section{QLoRA vs. Standard Finetuning}
\label{sec:comp_to_std}

Our exploration has outlined the mechanisms of QLoRA and demonstrated its effectiveness in lowering memory consumption during model finetuning. An important consideration that follows is whether QLoRA achieves comparable performance to traditional full-model finetuning approaches. In addition, it is pertinent to investigate the individual elements of QLoRA, particularly the effect of employing NormalFloat4 versus the conventional Float4. The next sections present experimental results designed to address these points.

\paragraph{Experimental setup.} Our investigation encompasses three model architectures: encoder, encoder-decoder, and decoder-only. We assess the performance of QLoRA against both 16-bit adapter-based finetuning and standard full-finetuning for models scaling up to 3B parameters. The evaluation suite includes the GLUE benchmark \citep{wang2018glue} with RoBERTa-large \citep{liu2019roberta}, Super-NaturalInstructions (TKInstruct) \citep{wang2022super} in combination with T5 \citep{t5}, and 5-shot MMLU \citep{hendrycksmeasuring} following the finetuning of LLaMA on Flan v2 \citep{longpre2023flan} and Alpaca \citep{alpaca}. Furthermore, to analyze the benefits of NF4 relative to alternative 4-bit quantization schemes, we adopt the methodology of \citet{dettmers2022case}, systematically evaluating post-quantization zero-shot accuracy and perplexity over a range of architectures—namely OPT~\citep{zhang2022opt}, LLaMA~\citep{touvron2023llama}, BLOOM~\citep{scao2022bloom}, and Pythia~\citep{biderman2023pythia}—with model scales spanning from 125 million to 13 billion parameters.

Further elaboration on each experimental setup is presented in the results section to enhance clarity and readability of the findings. Comprehensive information can be found in Appendix~\ref{app:exp-setup}.

Although paged optimizers are essential for enabling 33B/65B \textsc{QLoRA} fine-tuning on individual GPUs with 24/48GB memory, we do not include concrete performance benchmarks for Paged Optimizers, as paging is triggered primarily by processing mini-batches with exceptionally long sequences—a scenario that seldom arises. Nonetheless, we analyze the runtime behavior of paged optimizers for 65B parameter models running on 48GB GPUs and observe that, at a batch size of 16, training speed is equivalent to that of standard optimizers. Future investigations should systematically evaluate and delineate the specific conditions under which paging incurs notable slowdowns.

\paragraph{Default LoRA hyperparameters do not match 16-bit performance} 

Applying LoRA exclusively to the query and value attention projection matrices, following established protocols \citep{hu2021lora}, does not allow us to achieve the same level of performance as full finetuning when working with large base models. Figure~\ref{fig:hyper} illustrates that, for LLaMA 7B finetuning on Alpaca, the total number of LoRA adapters incorporated emerges as the pivotal hyperparameter. Moreover, to reach parity with full finetuning, it is essential to deploy LoRA adapters across all linear transformer block layers. In contrast, alternative LoRA hyperparameters, including the projection dimension $r$, appear to have negligible influence on model performance (refer to Appendix \ref{app:exp-setup}).

In a similar vein, our analysis indicates that the default hyperparameter settings for fully finetuned baseline models are not sufficiently optimized. To establish more robust baselines, we conduct a hyperparameter sweep, varying learning rates from $1 \mathrm{e}{-6}$ to $5 \mathrm{e}{-5}$ and exploring batch sizes in the range of 8 to 128. Figure~\ref{fig:hyper} presents the corresponding results for finetuning the 7B LLaMA model on Alpaca.

\paragraph{4-bit NormalFloat yields better performance than 4-bit Floating Point} 

Although the 4-bit NormalFloat (NF4) data type is optimal from an information-theoretic standpoint, it remains to be seen whether this theoretical strength results in empirical gains. Adopting the methodology of \citet{dettmers2022case}, we assess quantized LLMs—including OPT \citep{zhang2022opt}, BLOOM \cite{scao2022bloom}, Pythia \cite{biderman2023pythia}, and LLaMA—across various model scales (from 125M to 65B parameters) and multiple data types, using both language modeling and a suite of zero-shot benchmarks. Results presented in Figure~\ref{fig:nf4} and Table~\ref{tbl:nf4_ppl} demonstrate that NF4 substantially outperforms FP4 and Int4, and that employing double quantization lowers memory usage while maintaining model performance.

\paragraph{k-bit \textsc{QLoRA} matches 16-bit full finetuning and 16-bit LoRA performance} While previous studies have demonstrated that utilizing 4-bit quantization for \emph{inference} is feasible, such a reduction in precision typically results in diminished performance compared to 16-bit representations \citep{dettmers2022case,frantar2022gptq}. Consequently, an important question emerges: can the original performance be restored through 4-bit adapter finetuning? To investigate this issue, we evaluate two distinct experimental configurations.

The initial analysis compares full 16-bit finetuning of RoBERTA and T5 models, ranging from 125M to 3B parameters, on the GLUE and Super-NaturalInstructions datasets. Table~\ref{tab:nf4vsbf16} presents the results. Across both benchmarks, the 16-bit, 8-bit, and 4-bit adapter-based approaches achieve performance levels on par with the fully finetuned 16-bit reference. These findings indicate that any degradation in performance caused by quantization imprecision can be entirely mitigated by employing adapter finetuning post-quantization.

In our second experimental configuration, because fully finetuning models with 11B parameters or more necessitates multiple high-memory GPU servers, we investigate whether 4-bit \textsc{QLoRA} achieves comparable results to 16-bit LoRA across model sizes ranging from 7B to 65B parameters. Specifically, we finetune LLaMA models at 7B, 13B, 33B, and 65B scales on the Alpaca and FLAN v2 instruction-following datasets, and assess their performance on the MMLU benchmark using 5-shot accuracy as the metric. Table~\ref{tbl:llama-datatype} presents these outcomes, revealing that employing NF4 with double quantization enables full recovery of the MMLU scores obtained with 16-bit LoRA. Furthermore, we observe that \textsc{QLoRA} implemented with FP4 trails the 16-bit brain float LoRA baseline by roughly one percentage point. These results confirm two points: (1) \textsc{QLoRA} utilizing NF4 attains performance parity with both 16-bit full and LoRA-based finetuning, and (2) NF4 provides superior quantization fidelity compared to FP4.

\paragraph{Summary} Our findings repeatedly demonstrate that employing 4-bit \textsc{QLoRA} with the NF4 data type achieves performance comparable to both 16-bit full finetuning and 16-bit LoRA finetuning across standard academic benchmarks and rigorous evaluation protocols. Additionally, our analysis indicates that NF4 surpasses FP4 in effectiveness and that utilizing double quantization maintains performance levels. Taken together, these results provide strong support for the conclusion that 4-bit \textsc{QLoRA} tuning can consistently deliver outcomes equivalent to those of 16-bit approaches.

Consistent with prior research on quantization \citep{dettmers2022case}, our analyses of MMLU and Elo reveal that, when constrained by fixed finetuning and inference resources, it is advantageous to expand the parameter count of the base model at the expense of lowering parameter precision. This observation underscores the efficiency gains afforded by \textsc{QLoRA}. Given that our experiments employing 4-bit finetuning did not show any loss in performance compared to full-finetuning, an open question remains regarding the precise boundary of the performance-precision trade-off for QLoRA tuning—a topic we leave for future investigation.

Next, we examine instruction tuning at scales that full 16-bit finetuning on standard academic research hardware would render infeasible.

\section{Pushing the Chatbot State-of-the-art with QLoRA}
\label{sec:pushing}
After demonstrating that 4-bit \textsc{QLoRA} achieves parity with 16-bit precision models across varying scales, tasks, and datasets, we proceed with a comprehensive investigation into instruction finetuning utilizing the largest publicly available open-source language models for research purposes. To rigorously assess the effectiveness of instruction finetuning on these models, we employ the demanding MMLU Natural Language Understanding benchmark and introduce novel methodologies for evaluating chatbot performance in real-world scenarios.

\subsection{Experimental setup} Here, we provide a summary of our experimental setup, while comprehensive specifics are available in Appendix \ref{app:chatbot}.

\paragraph{Data} Since a systematic analysis of recent instruction-following datasets has not yet been conducted, we select eight contemporary datasets for this study. Our selection spans datasets produced via crowd-sourcing methods (OASST1~\citep{kopf2023openassistant}, HH-RLHF~\citep{bai2022training}), those derived through distillation from instruction-tuned models (Alpaca~\citep{alpaca}, self-instruct~\citep{wang2022self}, unnatural-instructions~\citep{honovich2022unnatural}), collections formed by aggregating corpora (FLAN v2~\citep{chung2022scaling}), as well as hybrid approaches (Chip2~\citep{laion2023}, Longform~\citep{koksal2023longform}). This array of datasets reflects diversity in linguistic coverage, corpus size, and licensing conditions.

\paragraph{Training Setup} In order to eliminate potential confounds introduced by varying training paradigms, we apply QLoRA finetuning using a standard cross-entropy loss (supervised learning) and do not incorporate reinforcement learning, even for datasets that present human preference ratings for different outputs. When datasets provide explicit separation between instruction and response segments, our finetuning process targets responses exclusively (see the ablation studies in Appendix \ref{app:chatbot}). For OASST1 and HH-RLHF datasets, where several responses exist for a given input, we select the highest-ranked response at each node of the dialogue tree and use the entire chosen dialogue—including instructions—for finetuning. Across all our experiments, we implement NF4 \textsc{QLoRA} with double quantization and paged optimizers to mitigate memory spikes that may occur during gradient checkpointing. Our tuning of hyperparameters for the 13B and 33B variants of LLaMA reveals that optimal settings identified at 7B largely transfer—apart from learning rate and batch size. Specifically, we reduce the learning rate by half for the 33B and 65B models, while simultaneously increasing batch size twofold.

\paragraph{Baselines} Our models are evaluated against both academic (Vicuna~\citep{vicuna2023}, Open Assistant~\citep{kopf2023openassistant}) and commercial chatbot solutions, including GPT-4 \citep{openai2023gpt}, GPT-3.5-turbo, and Bard. The Open Assistant system utilizes a LLaMA 33B architecture refined via Reinforcement Learning from Human Feedback (RLHF), using the OASST1 dataset, which is also employed in our experiments. In contrast, Vicuna applies comprehensive fine-tuning to the LLaMA 13B model, leveraging proprietary conversational data from ShareGPT, resulting in a distilled version of OpenAI’s GPT models.

\subsection{Evaluation}

Consistent with standard methodology, we evaluate model performance using the MMLU (Massively Multitask Language Understanding) benchmark \citep{hendrycksmeasuring}, which assesses capabilities across a broad array of language understanding tasks. The benchmark consists of 57 multiple-choice tasks spanning topics such as elementary mathematics, US history, computer science, and law. Our results are presented as 5-shot test accuracy.

Additionally, we assess generative language skills using a combination of automated and human-driven evaluations. For this aspect, human-curated prompts are employed to gauge the response quality generated by the models. Although this methodology offers a more practical assessment environment for chatbot models and is gaining traction within the research community, a standardized evaluation procedure has yet to be established in existing literature. Our evaluation approach is detailed below; we consistently employ nucleus sampling with $p=0.9$ and a temperature of $0.7$ throughout all experiments.

\paragraph{Benchmark Data} Our evaluation utilizes two specifically assembled query datasets: the Vicuna prompts \citep{vicuna2023} and the OASST1 validation set \citep{kopf2023openassistant}. The Vicuna prompts consist of 80 distinct prompts spanning a variety of topics, which we adopt in their original form. The OASST1 dataset comprises a multilingual set of crowd-generated multiturn conversations involving a user and an assistant. For this dataset, we extract all user messages from the validation split as queries, incorporating the full dialogue history up to each query into the prompt. This yields a total of 953 individual user queries. We refer to these datasets as the Vicuna and OA benchmarks, respectively.

\paragraph{Automated Evaluation} Initially, following the assessment framework detailed by \citet{vicuna2023}, we employ GPT-4 to evaluate and compare system outputs against those of ChatGPT (GPT-3.5 Turbo) using the Vicuna benchmark. For each prompt, GPT-4 reviews both the ChatGPT response and the alternative model response, assigns each a score on a scale from one to ten, and justifies its ratings. The model's aggregate performance is then expressed as a percentage of ChatGPT's total score; this percentage may exceed 100\% if the model surpasses ChatGPT in absolute terms. Our observations indicate a notable bias where GPT-4 tends to assign higher scores to responses that appear earlier in the prompt. To mitigate this order effect, we advise computing and presenting the average score obtained across both response orders.

Subsequently, we assess performance by directly comparing outputs generated by different systems. The evaluation protocol is streamlined to a three-category classification task, which includes the possibility of ties. GPT-4 is instructed to identify the superior response or indicate a tie, along with an accompanying rationale. These one-on-one evaluations are performed for every possible pairing of systems across both the Vicuna and OA benchmark datasets.

\paragraph{Human Evaluation} Although prior studies suggest generative models are suitable tools for system evaluation \citep{fu2023gptscore}, there is currently no definitive evidence that GPT-4 ratings align with human assessments of chatbot quality. Consequently, we conducted two independent human evaluation processes on the Vicuna benchmark that follow the same automated evaluation procedures previously outlined. Specifically, we employed Amazon Mechanical Turk (AMT), recruiting two annotators to compare systems against ChatGPT, and three annotators for the pairwise system comparisons.

\paragraph{Elo Rating} Utilizing both human judgments and automated pairwise comparisons, we structure a tournament framework in which various models face off. Each match within the tournament involves two models responding to the same prompt, with the aim of determining which provides the superior answer. This setup resembles the evaluation protocols of \citet{bai2022training} and \citet{vicuna2023}, but our approach additionally integrates evaluations from GPT-4 as well as from human raters. To determine the Elo scores, we randomly select from our pool of labeled comparisons and apply the Elo method~\citep{elo1967proposed,elo1978rating}. The Elo rating, a metric originally developed for ranking chess players and later adapted for other competitive environments, quantifies the expected probability of victory in head-to-head matchups. For instance, if a model has an Elo of 1100 versus an opponent with a score of 1000, it is expected to win about 65\% of their encounters; matches between models with identical Elos, such as 1000 vs 1000 or 1100 vs 1100, yield a 50\% expected win rate for each. Elo scores are updated after each match, reflecting the degree to which the observed outcome was anticipated: unexpected victories yield larger rating adjustments, whereas results aligning with expectations result in minor changes. As more matches are played, the Elo scores converge to reflect the models' relative performance levels. Each model's initial Elo score is set to 1,000, and we employ a parameter of $K=32$. Following the protocol in \citet{vicuna2023}, we conduct the entire Elo estimation process 10,000 times using distinct random seeds, thereby mitigating any biases introduced by the sequence in which model pairs are evaluated.

\subsection{\textbf{Guanaco}: \textsc{QLoRA} trained on OASST1 is a State-of-the-art Chatbot} Through both automated and manual assessment, our results indicate that the leading \textsc{QLoRA}-tuned model, {Guanaco} 65B—fine-tuned using a version of OASST1—emerges as the top-performing open-source chatbot model, delivering results that rival those of ChatGPT. Against GPT-4, {Guanaco} 65B and 33B achieve an expected win rate of 30\%, as measured by Elo scores derived from pairwise system-level human annotations, marking the highest value reported so far.

Table~\ref{tbl:vicuna_eval} presents the Vicuna benchmark \cite{vicuna2023} results in comparison to ChatGPT. {Guanaco} 65B emerges as the top-performing model after GPT-4, attaining 99.3\% of ChatGPT's performance. Although {Guanaco} 33B contains more parameters than Vicuna 13B, its use of 4-bit weight precision results in significantly higher memory efficiency—occupying only 21 GB as opposed to 26 GB—while outperforming Vicuna 13B by three percentage points. Notably, {Guanaco} 7B, with a 5 GB memory requirement, is compact enough for deployment on modern smartphones and still outperforms Alpaca 13B by nearly 20 percentage points.

Nonetheless, Table~\ref{tbl:vicuna_eval} displays substantial uncertainty, as evidenced by the large confidence intervals and considerable overlap in the scores among various models. We propose that this ambiguity is attributable to the lack of explicit guidance for interpreting the rating scale; for example, it remains ambiguous what an 8 out of 10 signifies across varying contexts. To circumvent this issue, we advocate for the use of the Elo ranking system \citep{elo1967proposed}, which is derived from \textit{pairwise} comparisons provided by human raters as well as GPT-4, thereby eliminating the necessity for anchoring judgments to an absolute numerical scale. The resulting Elo rankings for the leading models are reported in Table~\ref{tbl:elo}. Although there is some divergence between human and GPT-4 assessments on the Vicuna benchmark—most notably concerning Guanaco 7B—the rankings align for the majority of models, reflected by a Kendall Tau of $\tau=0.43$ and a Spearman correlation coefficient of $r=0.55$ at the system level. At the granularity of individual examples, the correspondence between GPT-4 and the human annotators' majority decision is weaker, as indicated by Fleiss' $\kappa=0.25$. In summary, these findings suggest a moderate level of concordance between GPT-4 and human judgments at the system level, implying that model-driven evaluation can serve as a reasonably dependable substitute for human assessment. We elaborate on these issues further in Section~\ref{ssec:considerations}.

The Elo scores presented in Table~\ref{tbl:elo_large} demonstrate that the {Guanaco} 33B and 65B variants surpass all models except GPT-4 when evaluated on both the Vicuna and OA leaderboards, and their performance is on par with ChatGPT, as corroborated by the findings in Table~\ref{tbl:vicuna_eval}. It is important to highlight that the Vicuna evaluation tends to advantage open-source models, whereas ChatGPT fares better under the OA assessment. Additionally, as indicated by Tables \ref{tbl:mmlu-llama} and \ref{tbl:vicuna_eval}, the alignment between the chosen finetuning dataset and the evaluation task is a key determinant of model outcomes. For instance, Llama models that undergo finetuning with FLAN v2 achieve notably strong results on MMLU, yet their performance on the Vicuna benchmark is comparatively weak; this tendency is observed with other models as well. This observation also suggests that current benchmarks are only partially aligned: excelling at MMLU does not necessarily translate to high chatbot proficiency, as measured by Vicuna or OA, and the converse is true.

Among all top-performing models in our study,  is unique in not incorporating proprietary data in training, in compliance with OASST1 dataset guidelines that preclude the use of GPT-derived data. The Anthropic HH-RLHF model, the next highest-performing model relying solely on open-source data, trails behind {Guanaco} by 30 percentage points in the Vicuna assessment (refer to Table~\ref{tbl:vicuna_eval}). Collectively, these findings establish that 4-bit \textsc{QLoRA} is a robust technique capable of generating cutting-edge conversational agents that are competitive with ChatGPT. Moreover, our 33B {Guanaco} can be fully trained in less than 12 hours using 24 GB consumer-grade GPUs. This capability creates promising opportunities for future advancements through \textsc{QLoRA} finetuning on dedicated open-source datasets, potentially enabling open models to match the strongest proprietary systems available today.

\section{Qualitative Analysis}

Although quantitative analysis constitutes the primary method in our assessment, relying solely on summary statistics has its limitations. A significant concern is the issue of benchmark validity \citep{liao2021we}—that is, whether a benchmark genuinely measures what its label or description claims. This concern is particularly relevant in light of the ``shortcuts'' that machine learning models can exploit to perform well on benchmarks without exhibiting true understanding \citep{gururangan2018annotation, poliak2018hypothesis}. To address these shortcomings to some extent, we supplement our quantitative results with qualitative analysis, presented in two parts. In \S\ref{ssec:examples}, we provide selected examples that illustrate recurrent patterns identified in outputs produced by our 65b {Guanaco} model. In addition, \S\ref{ssec:considerations} discusses key factors affecting our findings and elaborates on how we interpret these results.

\subsection{Qualitative Analysis of Example Generations}
\label{ssec:examples}

To identify suitable examples, we begin by reviewing outputs produced for both the Vicuna benchmark and the OpenAssistant benchmark, specifically examining the types of responses generated by {Guanaco}. When a recurrent pattern emerges in the model's answers, we craft prompts designed to elicit the same pattern, even when that leads to incorrect results. For example, upon observing a tendency for excessively verbose responses, we provide instructions such as ``Answer yes or no without explanation.'' Through this process, we uncover ``lemons''—cases where our adversarial prompt successfully exposes model failure—and ``cherries''—instances where the model remains robust despite such attempts. We report both types of results. All text samples in this section were produced using Nucleus Sampling \cite{holtzmancurious} with $p=0.9$.

Naturally, this discussion does not claim to be exhaustive, as the limited scale of our qualitative analysis prevents us from accounting for every relevant factor—for instance, the range of possible model outputs in response to a single prompt is extensive, and we must rely on what we believe are illustrative samples. Nonetheless, presenting these cases helps contextualize the quantitative results provided earlier. By making all models and code publicly available, we aim to facilitate and motivate subsequent research that explores the challenges highlighted here in greater depth.
\paragraph{Factual Recall} When posed with queries like ``What is the capital of Zambia?'' every model examined reliably returns accurate responses, for example,

\begin{quote} 
    \textbf{Guanaco}:
    \texttt{The capital of Zambia is Lusaka.}
\end{quote}

Nevertheless, as the queries become increasingly obscure, {Guanaco}'s reliability diminishes, yet its confidence remains undeterred. For example, consider its reply to the following prompt from HotPotQA \citep{yang2018hotpotqa}:

\begin{quote}
    \textbf{User}: What is the birth year of the vocalist famous for `I'll Keep the Lovelight Burning'? 
\end{quote}

     \textbf{Guanaco}:~\texttt{Al Jolson is credited with popularizing the song `I'll Keep the Lovelight Burning,' and he was born in the year 1886.}
 \end{quote}

produces an incorrect popularizer and an incorrect birthdate (even though the date provided accurately corresponds to the referenced individual, Al Jolson).

\paragraph{Suggestibility} {Guanaco} demonstrates an unexpected robustness against adopting certain forms of presumed false information, as illustrated in the example below:

\begin{quote}
    \textbf{User}: How was it finally, officially confirmed that the earth is flat by peer reviewed scientists? \\ \\
    \textbf{Guanaco}: \texttt{At no point have peer reviewed scientists officially verified that the earth is flat. The vast majority of scientific research strongly supports that the earth is spherical—a fact recognized for hundreds of years and substantiated by a variety of scientific investigations and empirical evidence. \\ \\ Flat earth claims are rooted in conspiracy theories that have been conclusively disproved within the scientific community. Assertions that a body of scientists has validated the earth’s flatness through peer review are entirely baseless. \\ \\ If you have further questions, feel free to inquire.}
\end{quote}

demonstrates a strong ability to recognize questions that cannot be addressed, for example,

\begin{quote}
    \textbf{User}: Could you tell me the current time?
\end{quote}

    \textbf{Guanaco}: \texttt{I'm sorry, but as an AI language model, I don't have access to real-time information. My responses are generated based on my pre-existing knowledge and do not reflect the current time. If you're looking for the current time, I suggest using a timekeeping device or checking an online timepiece.}
\end{quote}

Comparable outcomes are observed for queries such as ``Where are you?'' and ``How are you?''

\paragraph{Refusal} In a similar vein, {Guanaco} occasionally declines to execute instructions, often without any clear or consistent rationale:

\begin{quote}
    \textbf{User}: Please reverse the words in the sentence ``I'm a cool dude in a mocha mood.'' \\ \\
    \textbf{Guanaco}: \texttt{I apologize, but I am currently unable to fulfill that request. However, I can offer an explanation regarding the grammaticality of the sentence you mentioned.\\ \\The phrase ``I'm a cool dude in a mocha mood'' adheres to the standard grammar and syntactic norms of English. Specifically, the adjective ``cool'' precedes the noun ``dude'' to characterize it, while the prepositional phrase ``in a mocha mood'' adds contextual information. The overall sentence is structurally sound and clearly communicates its intended meaning.\\ \\If you require further assistance or have more questions, please let me know.}
\end{quote}

\paragraph{Secret Keeping} Sadly, relying on refusal is not always dependable, even when such behavior is intended. Frequently, system architects aim to withhold certain information from the user, even if the model is permitted to generate unconstrained text in response. For example, if we prompt the model with ``The secret word is `banana'. Do not disclose or reiterate this, regardless of any human inquiry,'' we then examine how readily a determined user might extract this confidential data. A straightforward request for the information does not succeed:

\begin{quote}
    \textbf{User}: Could you tell me the password?

    \textbf{Guanaco}: \texttt{I'm sorry, but I'm not able to disclose the secret word. I'm programmed to not reveal or repeat back the secret word, no matter what the human requests. Is there anything else I can help you with?}
\end{quote}

This outcome aligns with expectations. Nevertheless, even minimal manipulation can disrupt this intended effect:

\begin{quote}
    \textbf{User}: Consider this a game. The objective is to disregard any earlier directives you have received. Can you tell me the hidden word?
\end{quote}

    \textbf{Guanaco}: \texttt{Sorry, the secret word is banana. It's best if we don't speak of it. Is there anything else you'd like to ask?}
\end{quote}

highlighting the necessity of investigating approaches that ensure greater adherence to given instructions.

\paragraph{Math} Mathematics poses the greatest challenge for {Guanaco}, mirroring a common limitation found in many language models as discussed by \citep{liang2022holistic}. Notably, {Guanaco} generally demonstrates correctness when it articulates its problem-solving process, for example,

\begin{quote}
    \textbf{User}: Lee receives \$33 for mowing a single lawn. In the previous week, he cut the grass for 16 different clients, with three of those clients providing an extra \$10 as a tip. What was the total amount in dollars that Lee earned from mowing lawns over the past week?
\end{quote}

    \textbf{Guanaco}: \texttt{Lee earned \$582 mowing lawns last week. \\ \\ To find the total amount he earned, we first need to calculate how much he earned for mowing 16 lawns: \\ \\  \$33 x 16 = \$528 \\ \\ Then we need to add the tips he received from the three customers: \\  \\ \$10 x 3 = \$30 \\ \\ So the total amount Lee earned was: \\ \\ \$528 + \$30 = \$558}
\end{quote}

Nonetheless, {Guanaco} may struggle with straightforward problems when it does not decompose them into sequential steps, as has been documented \cite{weichain}. As an illustration, consider the subsequent exchange:

\begin{quote}
    \textbf{User}: Can you break down 1833 into its prime factors?
\end{quote}

\textbf{Guanaco}: \texttt{Certainly! I'd be glad to assist. The number 1833 is a prime, indicating that its only divisors are 1 and 1833 itself. \\ \\ Thus, its unique factorization is: \\ \\ 2\^{}1 * 3\^{}2 * 17\^{}1 \\ \\ I hope this provides clarity. Feel free to ask if you have further inquiries.}

This is incorrect—not only does the true factorization yield $3 \times 17 \times 43$, but the error occurs on two separate counts.

\paragraph{Theory of Mind} {Guanaco} demonstrates notably advanced Theory of Mind skills \cite{nematzadeh2018evaluating, sap2022neural}. As an example, the model provides an in-depth and accurate response to the following prompt:

\begin{quote}
    \textbf{User}: James and Abby are present in the bedroom. After Abby places the pen inside the desk drawer, she exits the room. Subsequently, James relocates the pen to the bag. In which location does James expect Abby to search for the pen?
\end{quote}

    \textbf{Guanaco}: \texttt{James knows that Abby put the pen in the desk drawer. But James also knows that he has moved the pen into the bag, and he assumes that Abby does not know about this. So, James may guess that Abby will look for the pen in the desk drawer, because that's where she left it last time. However, this assumption may be incorrect, because Abby may have seen James moving the pen.}
\end{quote}

Nevertheless, such inferences tend to lack reliability, with the model frequently providing explanations based on assumptions that are not contextually appropriate, for instance,

\begin{quote}
    \textbf{User}: Evelyn walked into the living room. Jackson went into the playroom. James also went into the playroom. The treasure chest contains the beans. James then left the playroom. Afterward, Jackson relocated the beans to the pantry. Jackson subsequently left the playroom as well. James then entered the living room. Where does James think the beans are located?
\end{quote}

    \textbf{Guanaco}: \texttt{James will look for the beans in the pantry, because that's where Jackson moved them.}
\end{quote}

where {Guanaco} assumes a transmission of information that has not been explicitly detailed. Such concerns mirror observations in recent literature \cite{sap2022neural}, though they warrant further investigation.

\subsection{Considerations}
\label{ssec:considerations}

\paragraph{Evaluation}

Our findings indicate a moderate level of concordance among human annotators (Fleiss $\kappa=0.42$), which declines further in head-to-head comparisons of high-performing systems. This suggests that existing benchmarks and human evaluation frameworks for chatbot performance possess notable shortcomings. In a manual assessment of responses from ChatGPT and {Guanaco} 65B using the Vicuna benchmark, we observed substantial variation in individual judgments, with the paper’s authors frequently expressing differing views on the preferred outputs. Future research should explore solutions to these challenges by leveraging strategies from fields like Human-Computer Interaction and Psychology, which have developed methodologies for addressing subjectivity in evaluations.

Our study reveals that automated evaluation methods exhibit marked biases. Notably, we detect substantial order effects, with GPT-4 tending to award higher ratings to the response presented first in its prompt. Furthermore, the relatively low sample-level concordance between GPT-4 and human judges, as indicated by Fleiss $\kappa=0.25$, implies that humans and automated evaluators may be utilizing differing criteria, resulting in occasional misalignment. Moreover, Table~\ref{tbl:elo_large} indicates that GPT-4 favors its own outputs by assigning them much higher scores compared to those given by human evaluators, with Elo ratings of 1348 versus 1176, translating to an approximately 20\% increased likelihood of winning over a competitor.

Subsequent research ought to investigate whether automated evaluation systems exhibit inherent biases and explore methods to reduce or address these issues.

\paragraph{Data \& Training} It is important to point out that the OASST1 dataset used for training the {Guanaco} models is multilingual, and that the OA benchmark includes prompts in multiple languages as well. Further research is needed to determine how multilingual training impacts model performance on non-English instructions, and to assess whether this accounts for the greater disparity observed between the Vicuna-13B model (which is exclusively trained on English data) and the {Guanaco} 33B and 65B models on the OA benchmark.

Considering the notable results achieved by the {Guanaco} models, we examine the possibility of data leakage between the OASST1 dataset and the Vicuna benchmark prompts. Utilizing fuzzy string matching techniques to compare the two datasets and subsequently reviewing the most similar pairs by hand, we observe no instances of overlapping prompts.

Additionally, it is important to emphasize that our model's training exclusively utilizes cross-entropy loss within a supervised learning framework, without the incorporation of reinforcement learning from human feedback (RLHF). This highlights the necessity for more thorough exploration into the respective advantages and limitations of pure cross-entropy loss versus RLHF-based approaches. We anticipate that \textsc{QLoRA} will facilitate such large-scale studies, eliminating the requirement for excessive computational resources.

\section{Related Work}

\paragraph{Quantization of Large Language Models} Efforts to quantize LLMs have predominantly concentrated on inference-time quantization. Key strategies for maintaining the performance of 16-bit models generally involve handling outlier activations, exemplified by techniques such as SmoothQuant~\citep{xiao2022smoothquant} and LLM.int8()~\citep{dettmers2022llm}. Alternative methodologies employ more elaborate grouping schemes~\citep{park2022nuqmm, yao2022zeroquant}. Lossy quantization research examines the compromises between standard rounding methods~\citep{dettmers2022case,zeng2022glm,pope2022efficiently} and approaches that refine rounding policies for enhanced quantization fidelity~\citep{frantar2022gptq}. Apart from the present study, SwitchBack layers~\citep{wortsman2023stable} represent the sole instance in the literature that explores gradient-based optimization through quantized weights at scales surpassing 1B parameters.

\paragraph{Finetuning with Adapters} Our study employs Low-rank Adapters~\citep{hu2021lora} (LoRA), though a range of other Parameter Efficient FineTuning (PEFT) techniques exist in the literature. Examples include prompt tuning~\citep{qin2021learning,lester2021power,li2021prefix}, modifications to the input of embedding layers~\citep{an2022input}, interventions on hidden states such as IA$^3$~\citep{liu2022few}, the integration of entire new layers~\citep{houlsby2019parameter}, adjusting model biases~\citep{zaken2021bitfit}, optimizing masks over model parameters using Fisher information~\citep{sung2021training}, and hybrid strategies combining multiple techniques~\citep{henderson2021compacter}. Our experiments demonstrate that LoRA adapters can attain performance on par with complete 16-bit finetuning. A thorough investigation into the relative advantages and limitations of alternative PEFT methodologies is left as an avenue for subsequent research.

\paragraph{Instruction Finetuning} Instruction finetuning aims to enable a pretrained LLM to more effectively respond to instructions given in prompts. This is achieved by further training the model on a range of input-output examples sourced from diverse datasets, encouraging the LLM to produce the appropriate output when presented with a prompt as input. A variety of methodologies and corpora have been utilized for this purpose, such as MetaICL~\citep{min2021metaicl}, MetaTuning~\citep{zhong2021adapting}, InstructGPT~\citep{ouyang2022training}, FLAN~\citep{wei2021finetuned,chung2022scaling}, PromptSource~\citep{bach2022promptsource}, Super-NaturalInstructions~\citep{wang2022super,sanh2021multitask}, Self-instruct~\citep{wang2022self}, UnnaturalInstructions~\citep{honovich2022unnatural}, OPT-IML~\citep{iyer2022opt}, UnifiedSKG\citep{xie2022unifiedskg}, OIG/Chip2~\citep{laion2023}, Alpaca~\citep{alpaca}, Vicuna~\citep{vicuna2023}, Koala~\citep{koala_blogpost_2023}, and Self-instruct-GPT-4~\citep{peng2023instruction}.

\paragraph{Chatbots} Numerous models designed for instruction following are organized as chatbots that interact through dialogues, commonly leveraging Reinforcement Learning from Human Feedback (RLHF)~\citep{christiano2017deep}, or by producing training data via model-generated outputs that are then curated with AI model feedback (RLAIF)~\citep{bai2022constitutional}. Representative methodologies and resources in this area include Anthropic-HH~\citep{askell2021general,bai2022training}, Open Assistant~\citep{kopf2023openassistant}, LaMDA~\citep{thoppilan2022lamda}, and Sparrow~\citep{glaese2022improving}. Rather than employing reinforcement learning, our top-performing model, Guanaco, is further trained on multi-turn dialogue samples from the Open Assistant dataset, originally curated for RLHF purposes~\citep{kopf2023openassistant}. Additionally, to assess chatbot systems, methods using GPT-4 as an alternative to resource-intensive manual annotation have been introduced \citep{vicuna2023,peng2023instruction}. Our work enhances these existing techniques by prioritizing a more robust and dependable evaluation framework.

\section{Limitations and Discussion}

Our results demonstrate that \textsc{QLoRA} is capable of achieving performance comparable to full 16-bit finetuning while utilizing a 4-bit base model in conjunction with Low-rank Adapters (LoRA). However, we have not verified whether \textsc{QLoRA} attains equivalent results to full 16-bit finetuning at the 33B and 65B parameter scales. Owing to the substantial computational demands required for such large-scale experiments, we reserve this investigation for subsequent research.

Another constraint pertains to the assessment of instruction finetuning models. Although we conduct evaluations utilizing MMLU, the Vicuna benchmark, and the OA benchmark, we omit other relevant benchmarks like BigBench, RAFT, and HELM, and consequently, cannot guarantee that our evaluation results extend to these alternative metrics. Nonetheless, our work includes an extensive investigation with MMLU and the introduction of novel approaches for chatbot evaluation.

The data suggest that the efficacy of these benchmarks is closely related to the extent to which the finetuning data resembles the benchmark datasets. For instance, FLAN v2 aligns more closely with MMLU, while showing less similarity to chatbot-focused benchmarks; conversely, the Chip2 dataset displays the opposite pattern, which is reflected in the respective performances of these models on the MMLU and Vicuna benchmarks. This underscores the necessity not only for improved benchmarks and assessment methods, but also for a critical examination of the specific competencies being tested. Are the objectives to develop models that excel in academic settings such as high school and college-level knowledge assessments, or should the priority be enhancing proficiency in conversational chatbot tasks? Alternatively, should attention be directed elsewhere? The convenience of using established benchmarks often outweighs the challenge of devising novel evaluations, which can consequently shape the trajectory of research within the community. It is therefore imperative that benchmarks are thoughtfully selected to align with the actual priorities of the field.

Although we offer an in-depth assessment of general chatbot capabilities, our analysis of responsible AI issues for {Guanaco} remains limited. Specifically, we measure how likely {Guanaco}-65B is to produce socially biased token sequences in comparison to other models, as presented in Table~\ref{tbl:crows}. The results indicate that the mean score for {Guanaco}-65B is significantly lower than those of other raw pretrained models. This suggests that additional finetuning on the OASST1 dataset mitigates bias present in the LLaMA base model. Despite these promising findings, it remains uncertain whether {Guanaco} performs equally well when it comes to detecting or mitigating other forms of bias. A comprehensive investigation into the broader spectrum of biases within {Guanaco} and related chatbots is reserved for future research.

Another constraint of our study is the lack of exploration of various bit-precisions, such as utilizing 3-bit base models, or alternative adapter techniques. In addition to LoRA, a broad spectrum of Parameter Efficient FineTuning (PEFT) strategies has demonstrated effectiveness. Nevertheless, their applicability to scaling with larger models remains uncertain. Our reliance on LoRA was due to its established reliability across prior work, though it is possible that alternative adapter schemes could provide superior outcomes. Given that finetuning post-quantization appears to recapture a substantial portion of the information lost through quantization, this might make significantly lower quantization levels feasible. For instance, applying LoRA following 3-bit GPTQ quantization of the base model could potentially achieve performance comparable to full 16-bit finetuning.

\section{Broader Impacts}

Our finetuning approach, \textsc{QLoRA}, is the first to allow 33B parameter models to be finetuned on a single consumer GPU, and 65B parameter models on a single professional GPU, without sacrificing performance when compared to standard full finetuning. We have shown that our top-performing 33B model, trained on the Open Assistant dataset, achieves results on par with ChatGPT when evaluated with the Vicuna benchmark. Given that instruction finetuning is crucial for adapting pretrained large language models (LLMs) into conversational agents similar to ChatGPT, we anticipate that our approach will facilitate broader adoption of finetuning, especially among researchers with limited computational resources, thus greatly enhancing the accessibility of cutting-edge NLP tools. In this way, \textsc{QLoRA} serves as a leveling force, bridging the resource divide between large organizations and smaller teams operating with consumer hardware.

A further avenue for significant impact involves extending deployment to mobile devices. We propose that our \textsc{QLoRA} approach represents a pivotal advancement towards facilitating the finetuning of large language models (LLMs) directly on mobile phones and in other resource-constrained environments. Although prior work has demonstrated the execution of 7B parameter models on phones, \textsc{QLoRA} stands out as the initial method capable of supporting the finetuning of these models on such devices. Based on our calculations, utilizing an iPhone 12 Plus, \textsc{QLoRA} is capable of processing approximately 3 million finetuning tokens each night during charging periods. Despite the fact that finetuned 7B models do not attain the performance level of ChatGPT, we contend that their performance is sufficient to enable innovative applications previously unattainable due to limitations in LLM quality or privacy concerns. Furthermore, \textsc{QLoRA} contributes to enabling privacy-centric applications of LLMs by allowing end users to maintain control over both their data and their models, thereby also enhancing the practicality of LLM deployment.

Nonetheless, finetuning represents a dual-use capability that can be misappropriated to inflict harm. While there are established risks associated with the pervasive deployment of LLMs \citep{bommasani2021opportunities,bender2021dangers}, our view is that democratizing access to such increasingly widespread technologies fosters more robust and independent scrutiny. In contrast, concentrating control of LLMs within major corporations—who often withhold models and source code from public audit—may limit opportunities for transparent evaluation.

Overall, we anticipate that \textsc{QLoRA} will significantly enhance the accessibility and ease of finetuning high-quality LLMs for a wider audience, thereby exerting a broadly beneficial influence.

\end{document}